El projecte ExoMars és una missió espacial amb la finalitat de buscar vida al planeta Mart. En una primera fase, el 2016, constava d'un satèl·lit, l'\textit{ExoMars Trace Gas Orbiter}, en òrbita circular al voltant de Mart a 400~km d'altura, i d'un mòdul de descens, l'\textit{Schiaparelli}, que havia d'aterrar a Mart.

Però quan el mòdul de descens estava a 3,7~km d'altura sobre Mart, pràcticament aturat, els sistemes automàtics van interpretar erròniament que ja havia arribat a la superfície. Van aturar els retrocoets i el mòdul es va desprendre del paracaigudes. Com a resultat, l'\textit{Schiaparelli} es va precipitar en caiguda lliure.

\figura[0.45]{fig1.png}

\begin{apartat}{1}
Calculeu el període de l'\textit{ExoMars Trace Gas Orbiter}.
\end{apartat}

\begin{apartat}{1}
Determineu el valor de l'acceleració de la gravetat a la superfície de Mart i la velocitat a la qual la nau va impactar a la superfície. (Considereu que la gravetat és constant durant la caiguda i la fricció amb l'atmosfera de Mart és negligible.)
\end{apartat}

\dades{%
  Massa de Mart $= 6,42\times 10^{23}\ \mathrm{kg}$.\\
  Radi de Mart $= 3,38\times 10^{6}\ \mathrm{m}$.\\
  $G = 6,67\times 10^{-11}\ \mathrm{N\,m^2\,kg^{-2}}$.}
